\section*{Limitations}
\label{sec:limitations}


\paragraph{English-only corpus.} 
\dolma was curated to contain English data. 
As tools for language identification may have false negatives, \dolma might contain a small percentage of non-English data.
Traces of non-English data are unlikely to lead to any meaningful downstream performance on non-English tasks for any model trained on \dolma. 
Thus, \dolma reinforces the expectation of English being the ``default'' language for NLP. 

\paragraph{Representativeness of sources in \dolma.} 
As mentioned in \S\ref{sec:desiderata}, it is impossible to curate a corpus that is representative of all language model data curation practices. 
Further, many open and close language models are trained on content that cannot be acquired or redistributed, and thus could not be included in \dolma. 

\paragraph{Single model configuration for ablations.} 
The experimental setup we use to validate our data curation pipeline only covers a subset of model types used to create language models. 
For example, while many language models are in the 7 billion to 70 billion parameters range, we train 1 billion parameter models;
further, we did not investigate the use of any alternative architectures to dense auto-regressive transformer models.
This choice was dictated by the need to efficiently iterate over many possible configurations, but it might result in design decisions that are not relevant at larger model sizes.
We expect downstream model developers to scrutinize \dolma before using it to train their language models, similar to the process we sketch in~\S\ref{sec:using}.

\paragraph{Limited tasks in evaluation suite.}
As detailed in \S\ref{sec:experimental-methodology}, we select tasks that have been used to evaluate previous base language models, and that are not present in our training data (i.e., \dolma is not contaminated against them). 
As such, we can only assess a subset of tasks language models are routinely used for. 
For example, the effect of adding code to pretraining data cannot be fully measured until models are able to generate executable code; such capability is typically observed only after models are finetuned to follow instructions~\citep{muennighoff2023octopack,zhuo2024astraios}.  


\paragraph{Manual inspection and evaluation of \dolma is infeasible.}
Given the corpus size, it is impossible to fully inspect \dolma to assess its content. 
While tools like WIMBD~\citep{wimbd} and Data Portraits~\citep{Marone2023DataPR} aid programmatic inspection of subsets of data, they cannot provide an assessment of all documents in a corpus. 
As such, we cannot fully describe the properties of \dolma in terms of data distribution, content quality, and potential harms due to the inclusion or exclusion of particular content.


\section*{Ethical Considerations}
\label{sec:ethics}




\paragraph{Minimize risk of harm to individuals during data curation.}

Curating a pretraining corpus may introduce risk to individuals, either by facilitating access to information that is present in the corpus, or by enabling training of harmful models
that disclose personal information~\citep{Carlini2020-jx} or produce toxic content~\citep{gehman-etal-2020-realtoxicityprompts,Ngo2021-fn}.
To minimize these risks while meeting our stated goals, we engaged with legal and ethics experts 
early in the project and evaluated data design decisions based on their feedback on a case-by-case basis.
Broadly, we follow accepted practices when available (\textit{e.g.}, masking of certain personal identifiable information), and take a measured approach when diverging opinions exist in the literature (\textit{e.g.}, most effective approach to identify and remove toxic content). 
Further, we will provide tools to request data removal\footnote{\label{foot:removal}Available at {\href{https://docs.google.com/forms/d/e/1FAIpQLSfL6KzFR7xNJj6MPyV1uikIpj-VmrftC9mjty2nXzSClU2rnw/viewform}{\path{forms.gle/FzpUXLJhE57JLJ3f8}}}}
We believe in compromising on desired research artifact properties like model reproducibility, performance, and extensibility in cases of significant harm to individuals.

Besides a risk-based approach, alternative frameworks for considering the ethical implications of language model data have also been proposed.
Data stewardship~\citep{Jernite2022DataGI} seeks to create a framework to collect and reflect explicit interests of data owners. 
Data trusts~\citep{datatrust} or data licensing~\citep{licensing} can also enable explicit consent in sharing data for AI training. 
As no current state-of-the-art model is trained on data collected through these frameworks, these approaches would limit the representativeness goal stated in \S\ref{sec:desiderata}.
As these principles are adopted, we will consider them for future versions of \dolma.




\paragraph{Copyright and fair use considerations.} 
At the time of writing, the landscape governing applicability of copyright law and fair use doctrine (also known as ``fair dealing'') and language models is largely undetermined~\citep{Cooper2023GenerativeAILaw,Lee2023Talking}. 
In the United States, legal scholars and practitioners have suggested that training models on copyright content might constitute fair use~\citep{iSchoolUIUC2023,UCBerkeley2023,Henderson2023FoundationMA}, while also recognizing limitations of existing doctrine in this application~\citep{uspto-copyright-ai2}.
Further, legal assessments regarding the use of copyrighted data in language models vary widely depending on jurisdiction: 
in early 2024, Israel~\citep{Israel2022} and Japan~\citep{Technomancersai2023-rx} allow copyrighted content to be used for AI training data, although the latter is currently re-considering this framework.
While most datasets we used were curated with copyright and licensing in mind (e.g., open access papers in peS2o~\citep{peS2o}, open source repositories in the Stack~\citep{kocetkov2022stack}) or were already permissively licensed (e.g., Wikipedia is released under a Creative Commons license), we recognize that large web crawls may also contain copyrighted material. 
Yet, given current tools, it's not possible to reliably or scalably detect copyrighted materials in a corpus of this size.
Our decision to curate and distribute \dolma factors in several considerations, including that all our data sources were publicly available and already being used in large-scale language model pretraining (both open and closed). 
We recognize that the legal landscape of AI is changing rapidly, especially as it pertains to use of copyrighted materials for training models. 


